%Ejercicio 1

\ejercicio{resuelto}{profesor}{
    Determine la solución general del sistema $u' = Au$ (tanto compleja como real) para cada una de las siguientes matrices reales $A$:
    \[
    \begin{pmatrix} -1 & 2 \\ -7 & 8 \end{pmatrix}, \quad 
    \begin{pmatrix} 1 & -1 \\ 5 & -3 \end{pmatrix}, \quad 
    \begin{pmatrix} 0 & 1 & 1 \\ 1 & 0 & 1 \\ 1 & 1 & 0 \end{pmatrix}, \quad 
    \begin{pmatrix} 10 & 4 & 13 \\ 5 & 3 & 7 \\ -9 & -4 & -12 \end{pmatrix}.
    \]
}{
    Para resolver cada sistema lineal homogéneo de coeficientes constantes $u' = Au$, procedemos a determinar los autovalores y autovectores (propios y generalizados, si fuese necesario) de cada matriz $A$.
    
    \begin{itemize}
        \item \textbf{Primera Matriz:} $A_1 = \begin{pmatrix} -1 & 2 \\ -7 & 8 \end{pmatrix}$
    
        El polinomio característico es $p(\lambda) = \det(A_1 - \lambda I) = \lambda^2 - \text{tr}(A_1)\lambda + \det(A_1) = \lambda^2 - 7\lambda + 6$.
        Sus raíces son los autovalores $\lambda_1 = 1$ y $\lambda_2 = 6$.
        
        Para $\lambda_1 = 1$:
        \[
            (A_1 - I)v_1 = \begin{pmatrix} -2 & 2 \\ -7 & 7 \end{pmatrix} \begin{pmatrix} x \\ y \end{pmatrix} = \begin{pmatrix} 0 \\ 0 \end{pmatrix} \implies v_1 = \begin{pmatrix} 1 \\ 1 \end{pmatrix}.
        \]
        
        Para $\lambda_2 = 6$:
        \[
            (A_1 - 6I)v_2 = \begin{pmatrix} -7 & 2 \\ -7 & 2 \end{pmatrix} \begin{pmatrix} x \\ y \end{pmatrix} = \begin{pmatrix} 0 \\ 0 \end{pmatrix} \implies v_2 = \begin{pmatrix} 2 \\ 7 \end{pmatrix}.
        \]
        
        Dado que ambos autovalores son reales y distintos, las soluciones reales y complejas tienen exactamente la misma forma algebraica (difieren únicamente en el cuerpo al que pertenecen las constantes $C_1, C_2$). La \textbf{solución general} es:
        \[
            u(t) = C_1 e^t \begin{pmatrix} 1 \\ 1 \end{pmatrix} + C_2 e^{6t} \begin{pmatrix} 2 \\ 7 \end{pmatrix}, \quad C_1, C_2 \in \mathbb{K} \text{ (donde } \mathbb{K}=\mathbb{R} \text{ o } \mathbb{C}).
        \]
        \item \textbf{Segunda Matriz:} $A_2 = \begin{pmatrix} 1 & -1 \\ 5 & -3 \end{pmatrix}$
    
        El polinomio característico es $p(\lambda) = \lambda^2 - (-2)\lambda + 2 = \lambda^2 + 2\lambda + 2$. Los autovalores son $\lambda = -1 \pm i$.
        Buscamos el autovector complejo asociado a $\lambda_1 = -1 + i$:
        \[
            (A_2 - (-1+i)I)v = \begin{pmatrix} 2-i & -1 \\ 5 & -2-i \end{pmatrix} \begin{pmatrix} x \\ y \end{pmatrix} = \begin{pmatrix} 0 \\ 0 \end{pmatrix} \implies (2-i)x - y = 0 \implies v_1 = \begin{pmatrix} 1 \\ 2-i \end{pmatrix}.
        \]
        
        La \textbf{solución general compleja} es una combinación lineal de las soluciones conjugadas:
        \[
            u_\mathbb{C}(t) = C_1 e^{(-1+i)t} \begin{pmatrix} 1 \\ 2-i \end{pmatrix} + C_2 e^{(-1-i)t} \begin{pmatrix} 1 \\ 2+i \end{pmatrix}, \quad C_1, C_2 \in \mathbb{C}.
        \]
        
        Para obtener un sistema fundamental real, utilizamos la fórmula de Euler $e^{(-1+i)t} = e^{-t}(\cos t + i\sin t)$ y separamos la parte real y la imaginaria de una de las soluciones complejas:
        \[
            e^{-t}(\cos t + i\sin t) \left[ \begin{pmatrix} 1 \\ 2 \end{pmatrix} + i\begin{pmatrix} 0 \\ -1 \end{pmatrix} \right] = e^{-t} \begin{pmatrix} \cos t \\ 2\cos t + \sin t \end{pmatrix} + i e^{-t} \begin{pmatrix} \sin t \\ 2\sin t - \cos t \end{pmatrix}.
        \]
        Así, la \textbf{solución general real} es:
        \[
            u_\mathbb{R}(t) = K_1 e^{-t} \begin{pmatrix} \cos t \\ 2\cos t + \sin t \end{pmatrix} + K_2 e^{-t} \begin{pmatrix} \sin t \\ 2\sin t - \cos t \end{pmatrix}, \quad K_1, K_2 \in \mathbb{R}.
        \]
        Otra forma de resolver el sistema sería:
        \[
            e^{tA} = \begin{pmatrix}
                1 & 0 \\
                2 & -1
            \end{pmatrix} \begin{pmatrix}
                e^{-t} \cos t & e^{-t} \sin t \\
                -e^{-t} \sin t & e^{-t} \cos t
            \end{pmatrix} \begin{pmatrix}
                1 & 0 \\
                2 & -1
            \end{pmatrix}
        \]
        \[
            = \begin{pmatrix}
                1 & 0 \\
                2 & -1
            \end{pmatrix} \begin{pmatrix}
                2 e^{-t} \sin t + e^{-t} \cos t & - e^{-t} \sin t \\
                2 e^{-t} \cos t - e^{-t} \sin t & - e^{-t} \cos t
            \end{pmatrix} = \begin{pmatrix}
                2 e^{-t} \sin t + e^{-t} \cos t & - e^{-t} \sin t \\
                4 e^{-t} \sin t + 2 e^{-t} \cos t - 2 e^{-t} \cos t + e^{-t} \sin t & - 2 e^{-t} \sin t + e^{-t} \cos t
            \end{pmatrix}
        \]
        Para hallar la solución general, multiplicamos $e^{tA}$ por un vector constante $C = \begin{pmatrix} C_1 \\ C_2 \end{pmatrix}$.
        \item \textbf{Tercera Matriz:} $A_3 = \begin{pmatrix} 0 & 1 & 1 \\ 1 & 0 & 1 \\ 1 & 1 & 0 \end{pmatrix}$
    
        El polinomio característico es $p(\lambda) = -\lambda^3 + 3\lambda + 2 = -(\lambda - 2)(\lambda + 1)^2$. La matriz es simétrica, por lo que diagonaliza siempre sobre $\mathbb{R}$. Los autovalores son $\lambda_1 = 2$ (simple) y $\lambda_2 = -1$ (doble).
        
        Para $\lambda_1 = 2$:
        \[
            (A_3 - 2I)v_1 = 0 \implies \begin{pmatrix} -2 & 1 & 1 \\ 1 & -2 & 1 \\ 1 & 1 & -2 \end{pmatrix}\begin{pmatrix} x \\ y \\ z \end{pmatrix} = 0 \implies v_1 = \begin{pmatrix} 1 \\ 1 \\ 1 \end{pmatrix}.
        \]
        
        Para $\lambda_2 = -1$:
        \[
            (A_3 + I)v = 0 \implies \begin{pmatrix} 1 & 1 & 1 \\ 1 & 1 & 1 \\ 1 & 1 & 1 \end{pmatrix}\begin{pmatrix} x \\ y \\ z \end{pmatrix} = 0 \implies x + y + z = 0.
        \]
        Tomamos dos vectores linealmente independientes que cumplan la ecuación, por ejemplo, $v_2 = \begin{pmatrix} -1 \\ 1 \\ 0 \end{pmatrix}$ y $v_3 = \begin{pmatrix} -1 \\ 0 \\ 1 \end{pmatrix}$.
        
        Dado que todos los autovalores y autovectores son puramente reales, las soluciones complejas y reales son idénticas. La \textbf{solución general} es:
        \[
            u(t) = C_1 e^{2t} \begin{pmatrix} 1 \\ 1 \\ 1 \end{pmatrix} + C_2 e^{-t} \begin{pmatrix} -1 \\ 1 \\ 0 \end{pmatrix} + C_3 e^{-t} \begin{pmatrix} -1 \\ 0 \\ 1 \end{pmatrix}, \quad C_1, C_2, C_3 \in \mathbb{K}.
        \]
        \item \textbf{Cuarta Matriz:} $A_4 = \begin{pmatrix} 10 & 4 & 13 \\ 5 & 3 & 7 \\ -9 & -4 & -12 \end{pmatrix}$
    
        El polinomio característico es $p(\lambda) = -\lambda^3 + \lambda^2 + \lambda - 1 = -(\lambda-1)^2(\lambda+1)$.
        
        Para $\lambda_1 = -1$:
        \[
            (A_4 + I)v_1 = \begin{pmatrix} 11 & 4 & 13 \\ 5 & 4 & 7 \\ -9 & -4 & -11 \end{pmatrix}\begin{pmatrix} x \\ y \\ z \end{pmatrix} = 0 \implies v_1 = \begin{pmatrix} 2 \\ 1 \\ -2 \end{pmatrix}.
        \]
        
        Para $\lambda_2 = 1$ (doble):
        \[
            (A_4 - I)v_2 = \begin{pmatrix} 9 & 4 & 13 \\ 5 & 2 & 7 \\ -9 & -4 & -13 \end{pmatrix}\begin{pmatrix} x \\ y \\ z \end{pmatrix} = 0 \implies v_2 = \begin{pmatrix} 1 \\ 1 \\ -1 \end{pmatrix}.
        \]
        Dado que el rango de $A_4 - I$ es 2, el espacio propio tiene dimensión 1. Requerimos de un autovector generalizado $w$ que satisfaga $(A_4 - I)w = v_2$:
        \[
            \begin{pmatrix} 9 & 4 & 13 \\ 5 & 2 & 7 \\ -9 & -4 & -13 \end{pmatrix} \begin{pmatrix} x \\ y \\ z \end{pmatrix} = \begin{pmatrix} 1 \\ 1 \\ -1 \end{pmatrix} \implies w = \begin{pmatrix} 1 \\ -2 \\ 0 \end{pmatrix}.
        \]
        
        Como los autovalores y vectores son todos reales, ambas soluciones coindicen formalmente en estructura. La \textbf{solución general} se construye usando el vector propio y el generalizado:
        \[
            u(t) = C_1 e^{-t} \begin{pmatrix} 2 \\ 1 \\ -2 \end{pmatrix} + C_2 e^{t} \begin{pmatrix} 1 \\ 1 \\ -1 \end{pmatrix} + C_3 e^t \left[ \begin{pmatrix} 1 \\ -2 \\ 0 \end{pmatrix} + t \begin{pmatrix} 1 \\ 1 \\ -1 \end{pmatrix} \right], \quad C_1, C_2, C_3 \in \mathbb{K}.
        \]
    \end{itemize}    
}

% Ejercicio 2
\ejercicio{enunciado}{no}{
    Determine las formas canónicas compleja y real de la matriz
    \[
    A = \begin{pmatrix} 
    2 & 2 & -1 & -1 \\ 
    -1 & 3 & -1 & 1 \\ 
    1 & 1 & 1 & 0 \\ 
    0 & 1 & -1 & 2 
    \end{pmatrix},
    \]
    así como la matriz exponencial $e^{zA}$ a partir de la forma canónica real de $A$.
}{

}

% Ejercicio 3
\ejercicio{enunciado}{no}{

}{

}

% Ejercicio 4
\ejercicio{enunciado}{no}{

}{

}

% Ejercicio 5
\ejercicio{enunciado}{no}{

}{

}

% Ejercicio 6
\ejercicio{enunciado}{no}{

}{

}

% Ejercicio 7
\ejercicio{enunciado}{no}{

}{

}

% Ejercicio 8
\ejercicio{enunciado}{no}{

}{

}

% Ejercicio 9
\ejercicio{enunciado}{no}{

}{

}

% Ejercicio 10
\ejercicio{enunciado}{no}{

}{

}

% Ejercicio 11
\ejercicio{enunciado}{no}{

}{

}

% Ejercicio 12
\ejercicio{enunciado}{no}{

}{

}

% Ejercicio 13
\ejercicio{enunciado}{no}{

}{

}

% Ejercicio 14
\ejercicio{enunciado}{no}{

}{

}

% Ejercicio 15
\ejercicio{enunciado}{no}{

}{

}
